% !TEX root = ../main.tex
%% Devoir surveillé n° 5 — chapitres 14 à 16
%% Résultats vérifiés par verifications/eval_ds.py.

\sujetEval{DS5}{MPSI · Devoir surveillé}{Devoir surveillé \numero5}
  {Durée : \textbf{4 heures}\\Calculatrice interdite\\Chapitres 14 à 16}{3}
  {Exercice 1/Permutations et déterminants/6 pts,Exercice 2/Intégration/8 pts,%
   Exercice 3/Dénombrement/6 pts}
\consigne{Les trois exercices sont indépendants. Dans l'exercice 3, on
justifiera chaque dénombrement par une construction explicite.}

\begin{exoEval}{6 points}{Permutations et déterminants}
\begin{questions}
  \item Soit $\sigma\in S_8$ définie par
    $\sigma=\begin{pmatrix}1&2&3&4&5&6&7&8\\3&5&7&1&2&8&4&6\end{pmatrix}$.
    Décomposer $\sigma$ en cycles à supports disjoints ; en déduire sa
    signature et son ordre.
  \item Soient $a,b\in\R$ et $D_n$ le déterminant de la matrice de taille
    $n$ dont les coefficients diagonaux valent $a$ et les autres $b$.
    Montrer que $D_n=\bigl(a+(n-1)b\bigr)(a-b)^{n-1}$. \emph{Ajouter toutes
    les colonnes à la première.}
  \item En déduire que $J_n-I_n$ est inversible pour $n\geqslant2$, où
    $J_n$ est la matrice dont tous les coefficients valent $1$.
  \item Sans le développer, montrer que
    $\begin{vmatrix}1&1&1\\a&b&c\\b+c&c+a&a+b\end{vmatrix}=0$.
\end{questions}
\end{exoEval}

\begin{exoEval}{8 points}{Intégration}
Pour $n\in\N$, on pose $I_n=\displaystyle\int_0^{1}\frac{x^{n}}{1+x^{2}}\,\mathrm{d}x$.
\begin{questions}
  \item Calculer $I_0$ et $I_1$.
  \item Montrer que $(I_n)$ est décroissante, positive, et que
    $I_n\leqslant\dfrac{1}{n+1}$.
  \item Montrer que $I_n+I_{n+2}=\dfrac{1}{n+1}$, puis que
    $\dfrac{1}{2(n+1)}\leqslant I_n\leqslant\dfrac{1}{2(n-1)}$ pour
    $n\geqslant2$. En déduire un équivalent de $I_n$.
  \item Déterminer $\displaystyle\lim_{n\to+\infty}\sum_{k=1}^{n}\frac{1}{\sqrt{n^{2}+kn}}$.
  \item À l'aide de l'inégalité de Taylor-Lagrange, montrer que :
    \[\forall x\geqslant0,\qquad
      \Bigl|\ln(1+x)-x+\frac{x^{2}}{2}\Bigr|\leqslant\frac{x^{3}}{3}.\]
\end{questions}
\end{exoEval}

\begin{exoEval}{6 points}{Dénombrement}
\begin{questions}
  \item Combien d'anagrammes le mot MISSISSIPPI possède-t-il ?
  \item On tire $5$ cartes dans un jeu de $32$ ($8$ hauteurs, $4$
    couleurs). Combien de mains contiennent un \emph{full} (trois cartes
    d'une hauteur et deux d'une autre) ? une \emph{double paire} (deux
    paires de hauteurs distinctes et une cinquième carte d'une troisième
    hauteur) ?
  \item De combien de façons peut-on distribuer $10$ bonbons identiques à
    $4$ enfants, chacun en recevant au moins un ? Et sans cette
    contrainte ?
\end{questions}
\end{exoEval}
